% BIOLOGICAL METHODS AND EXACT CERTIFICATES
% Scientific reference text in TeX notation; this is not a standalone build.
% The seven complete tables are in biological_tables.tex. Commands and data
% locators are in biological_reproduction.txt. Bibliographic credits follow.

\section{Model identity and the control convention}
\label{bio:model-convention}
We evaluate four fully specified Boolean maps. Two code variants from the
regulatory-network benchmark collection have 17 and 29 nodes;
The 17-node and 29-node rule tables give their complete update
rules and targets. Two additional curated models describe CD4$^+$ T-cell
signaling and macrophage activation and have 188 and 321 coordinates,
respectively, including their external inputs. Their exact source identities,
input conditions, target construction and verification protocol are given
in the larger-model protocol section. Complete source rules, explicit input completion, ordered target
vectors and source checksums accompany the computational package.

The principal optimization uses synchronous updates; the fair-coordinate-update section treats two selected targets under coordinate-fair asynchronous updates. A chosen node is held permanently at its
specified target value, beginning at time zero. All other nodes follow
their Boolean rules, even when their current value happens to equal a
target value. All nodes are eligible, with unit intervention costs.
Success means eventual arrival at the complete fixed target from
\emph{every} initial state compatible with the chosen pins. No sampled
initial-state criterion, marginal entropy threshold, or automatic
freezing of deterministic coordinates is used as a success test.

We distinguish the total cost, including persistent external-input clamps,
from the additional cost when a specified environmental baseline is imposed
without cost. For the two smaller variants, the identity inputs are
$B_{17}=\{1,2,3\}$ and $B_{29}=\{1,9\}$. Any intervention that succeeds
from every compatible initial state must contain these inputs: without
its own pin, an identity coordinate initially different from target remains
different forever. Their respective total optima are $3,3$, and their
conditional optima are $0,1$. The same accounting is used for the larger
models, with the full external-input sets specified below.

\paragraph{Frozen sources.}
The repository is \url{https://github.com/ziyuxuan67-maker/QL-and-RTS},
at commit \texttt{29b4f8c5e8e15c71735ab643813cfdb6de7d8d77}.
The extraction reads the target literals and the polynomial assignments
in \path{17nodeQL.py} and \path{29nodeQL.py}. It does not execute their
learning, mean-field, automatic-freezing, plotting, or file-writing
wrappers. Every local Boolean corner is evaluated with exact integer
arithmetic. The 17-node and 29-node frozen source expressions have
$80$ and $88$ syntactic local corners, respectively; all are Boolean,
and both supplied targets are fixed points. The 29-node rule templates
in the source template collection agree with the frozen code at all
$88$ corners and have the same target. The template table has $82$
corners because its Fas row already omits two functionally irrelevant
inputs present in the frozen polynomial expression.

\section{Complete update tables}
\label{bio:update-tables}
For a row whose displayed ordered inputs are $(p_1,\ldots,p_{a_i})$,
the truth-word index is $k=\sum_{j=1}^{a_i}2^{j-1}x_{p_j}$.
Thus the first input supplies the least significant local bit, and the
first truth-word character corresponds to $k=0$. A repeated occurrence
of a variable always uses the same physical bit; it is not an independent
input. These tables, together with the target column, define the full
physical maps without biological-name ambiguity. The JSON and BNET
exports use neutral labels \texttt{n01}, \texttt{n02}, and so on.
\subsection{Segment-polarity code variant}
% Complete table tab:rules17: biological_tables.tex

\subsection{T--LGL code variant}
% Complete table tab:rules29: biological_tables.tex


\section{Certificate details and independent checks}
\label{bio:certificate-checks}
Let $y=x\oplus x^\star$. With a baseline $B$ fixed, exact coordinate-cube
support starts from $S_0=V\setminus B$ and deletes coordinates whose
Boolean outputs cannot deviate from target when all deviations outside
$S_t$ are zero. The greatest fixed set $R_B$ captures every controlled
trajectory. The main article proves that pinning optimization on the
full model is equivalent to pinning optimization on the physical
restriction to $R_B$, including for additional pins outside the residual.

\subsection{T--LGL copy certificate}
With $B=\{1,9\}$ the support cardinalities are
$(27,21,12,8,5)$ and
$R_B=\{15,16,17,18,19\}$. The restricted deviation rules are
\[
 y_{15}^+=y_{16},\quad y_{16}^+=y_{15},\quad
 y_{17}^+=y_{15},\quad y_{18}^+=y_{17},\quad y_{19}^+=y_{18}.
\]
The two-node cycle requires and permits exactly one additional pin.
The complete list of minimum total interventions is
$\{1,9,15\}$ and $\{1,9,16\}$.
Complementing the target at all five residual nodes produces a
competing fixed point when only $B$ is pinned. Both three-pin
interventions have empty support after four updates, with cardinality
sequences $(26,18,8,3,0)$ and $(26,19,9,4,0)$.

The four-update bound is attained. For either intervention, initialize
at the supplied target except for node $10$, which is changed from $0$
to $1$. Exact update then propagates a deviation through the downstream
chain before it vanishes at time four. The complete witness trajectories
and competing fixed point are recorded in
\path{results/tlgl_coded_29_certificate.json}. These are finite symbolic
certificates, not an enumeration of the $2^{26}$ compatible full states.
The cycle structure and controller $\{1,9,15\}$ are already reported
for a related model in Zhu et al.\ \cite[Example 6.2]{Zhu2024}.

\subsection{Segment-polarity nonlinear residual}
With $B=\{1,2,3\}$, support sizes are $(14,12,11,9,7)$ and
$R_B=\{4,5,10,11,13,16,17\}$. In that order, write the deviations
as $(a,b,c,d,e,f,h)$. Their exact restricted updates are
\[
 a^+=h\wedge(a\vee f),\quad b^+=a,\quad c^+=f\wedge h,\quad
 d^+=d\wedge\neg c,\quad e^+=d,\quad f^+=e,\quad h^+=e.
\]
The copy certificate fails, but the physical dynamics still converge.
The main article proves $d(t)=0$ for $t\ge4$ and complete residual
extinction by time $8$. The residual initial state $(0,0,0,1,0,0,0)$
attains this bound. Exhaustive enumeration checks all $128$ residual
states and all $16384$ compatible full-network states. The full model's
maximum transient is $11$, attained, for example, by the physical initial
state with integer encoding $99$ when node $1$ is the least significant
bit. The target has encoding $91139$. The trajectory is preserved in
\path{results/drosophila_coded_17_certificate.json}.

Only the three mandatory inputs are needed, so $\{1,2,3\}$ is the unique
minimum total intervention. This result explains why a sufficient
support-collapse search can return an unnecessary fourth pin. Pinning
node $11$ in addition to $B$ yields support sizes
$(13,10,8,5,3,2,1,0)$, but that extra pin is not required for physical
global convergence. We retain the smaller exact optimum.

\paragraph{Reachable-support curve.}
For the 17-node example, let $X_t$ be the image at time $t$ of the
complete initial cube compatible with $B$. The exact curve in the main
article is $|\bigcup_{x\in X_t}\supp(x\oplus x^\star)|$; the comparison
curve is the iterated coordinate-cube support $|S_t|$. At times $0$--$11$,
the exact values are $(14,12,11,9,7,7,7,6,5,3,1,0)$, while the cube
values are $(14,12,11,9,7,7,7,7,7,7,7,7)$. The respective numbers of
distinct reachable states are
$(16384,720,135,56,24,12,7,5,4,3,2,1)$.
All values are exact finite-state counts.

\section{Eight-model inventory and biological interpretation}
\label{bio:source-qualifications}
This inventory distinguishes reported intervention cardinalities from
independently verified optima across eight benchmark names. The certified
computations concern the two fully specified small implementations and the two
larger curated models described below. Source availability and complete-model
identity are recorded explicitly; missing definitions are not failed certificates.
The eight-model inventory table lists the source-reported cardinalities and
verification status.


% Complete table tab:inventory: biological_tables.tex


\paragraph{Segment-polarity rule version.}
\label{bio:cia-cir-comparison}
The code's node ordering corresponds to the supplied single-cell model
file associated with Parmer et al.\ \cite{Parmer2022}, with the source
label \texttt{nhhnHH} shortened to nHH in the source benchmark description.
The rules for CIA and CIR (nodes $16,17$) differ between those sources.
The archived code uses the dynamic SMO coordinate, while the comparison
file substitutes its expression in terms of PTC and the neighboring
hedgehog input. Fifteen rules agree, and two differ when compared over
all $96$ corners of the union input scopes. Substitution can change
synchronous transients, so the coded equations are retained. The indexed
certificate concerns these equations and their supplied target.

The comparison rules below are from
\path{biological/provenance/source_files/drosophila_single_cell.txt}
under \path{biological/provenance/source_files/}; the coded rules are the corresponding rows
of the 17-node rule table. Here CI, SMO, PTC and \texttt{nhhnHH}
are nodes $15,13,11,3$, respectively.
\begin{center}\small
\begin{tabular}{@{}lcc@{}}
\toprule
Node & Coded rule & Comparison rule\\
\midrule
CIA ($16$) & $\mathrm{CI}\wedge(\mathrm{SMO}\vee\mathrm{nhhnHH})$
 & $\mathrm{CI}\wedge(\neg\mathrm{PTC}\vee\mathrm{nhhnHH})$\\
CIR ($17$) & $\mathrm{CI}\wedge\neg\mathrm{SMO}\wedge\neg\mathrm{nhhnHH}$
 & $\mathrm{CI}\wedge\mathrm{PTC}\wedge\neg\mathrm{nhhnHH}$\\
\bottomrule
\end{tabular}
\end{center}

\paragraph{T--LGL labels and target.}
\label{bio:fas-variant}
The explicit node map of Zhu et al.\ \cite[Eq.~(17), p.~184]{Zhu2024}
identifies nodes $1,9,15,16,20,21$ as IL15, PDGF, SPHK, S1P,
Caspase and Apoptosis, respectively. The coded target values at those
nodes are $(1,0,1,1,0,0)$. Thus the complete coded target should not be
called apoptotic without a separate biological reconciliation.
The archive's selected set uses SPHK (node $15$), while the source benchmark description names S1P (node $16$); both alternatives are optimal for
the coded map. The code uses $F_{18}=\neg x_{17}$, whereas the cited
model has $F_{18}=(\neg x_{17})\vee[(\neg x_1)\wedge(\neg x_{11})]$.
The certificate preserves the coded rule. The prescribed baseline
$x_1=1$ makes this particular difference vanish in the controlled
comparison, but it does not establish equality of the unpinned models.

\paragraph{ABA implementation.}
The supplied 44-node polynomial for node $16$ expands to
$x_{13}x_{15}+x_{22}x_{23}$, which equals $2$ when all four inputs are
$1$. It is therefore not a Boolean rule. The supplied target is nevertheless
a fixed point of the polynomial assignments; target invariance alone does
not establish that all Boolean states have Boolean successors.
The four coded identity inputs are consistent with counting 40 dynamic
coordinates plus four external coordinates. This count by itself does not
resolve the rule or naming discrepancies. Clipping the polynomial or
replacing its sum by OR would change the implementation. An authoritative
index map and corrected rule set, together with target and update
conventions, are needed. The original published ABA model uses asynchronous
updates and includes oscillating components \cite{Li2006}; a full fixed-point
objective is not interchangeable with a stomatal-closure phenotype.

\paragraph{Remaining model inputs.}
For each of the five unavailable models, the required files are a complete
versioned rule export, its ordered node list, the full target or phenotype
predicate, and the external-input and update conventions. The cholesterol
row also needs reconciliation: it reports four intervention nodes but
lists only three names. The machine-readable inventory gives the exact
missing fields for every model. None is filled by guessing a different
public model version.

\section{Larger curated models and the fixed-input protocol}\label{sec:largeprotocol}
The larger benchmarks use the CD4$^+$ T-cell signaling model associated
with Conroy et al.\ \cite{Conroy2014} and the macrophage activation model
associated with Raza et al.\ \cite{Raza2008}. We take the exact Boolean
exports from the Biodivine Boolean Models (BBM) collection~\cite{Pastva2023},
identifiers 050 and 001, at repository commit
\texttt{314b2fed94689c38d3aaaf8101ed622b3501c7f7} of
\url{https://github.com/sybila/biodivine-boolean-models}.
The packaged source BNET files retain their original bytes. BBM records
154 regulated coordinates and 34 inputs for CD4 signaling, and 302
regulated coordinates and 19 inputs for macrophage activation. Thus the
completed synchronous maps have 188 and 321 coordinates.
Both source metadata records identify Cell Collective~\cite{Helikar2012}. The macrophage source is a later Cell Collective/BBM Boolean realization.
Its 321-coordinate count comes from the executable export; Raza et al.'s
original pathway diagram contains 295 nodes~\cite{Raza2008}.

BBM leaves the source-node update functions unspecified. We complete
exactly those omitted input rows as identity rules $F_i(x)=x_i$ and retain
every supplied regulated-node formula. Consequently, input values are
persistent once specified. This is an explicit mathematical convention
for the downloaded model versions. The metadata also records BBM's prior
curation: five CD4 regulation signs were marked unknown, and four incoming
macrophage regulations at IFNB1gene were marked non-essential to agree
with the update functions. These annotations are preserved, without any
further biological-rule modification.

\paragraph{Source identity and complete model data.}
The SHA256 digest of the downloaded CD4 BNET file is
\begin{center}\footnotesize\texttt{24722017d5be15ab31cf9a8c3dfc31139a80313c7a69196f8ab77b0945444209},\end{center}
and that of the macrophage BNET file is
\begin{center}\footnotesize\texttt{ea60f90ecab55c939f44648dc24608d5779916f7d1df54200f1e85681a4951d4}.\end{center}
The computational package contains these source files, their metadata,
immutable download URLs, the exact input-identity completion convention
and an explicit ordered list of every coordinate. Numeric state encodings always refer
to that recorded order; gene names are never matched by an inferred
abbreviation. The complete target vectors are distributed with the
case records, so reproducing a case requires no phenotype interpretation
or reconstruction of omitted values.

\subsection{Complete input conditions and target selection}
The protocol is frozen in \path{larger_networks/protocol/contexts.json},
with SHA256 digest
\begin{center}\footnotesize\texttt{f2a326e2a2201f011a0482db9c811393793691d46183bcd64bb1f67b1c5df41b}.\end{center}
Four all-zero/all-one endpoint conditions had already been inspected in
preliminary feasibility work and are retained explicitly as such. The
seven focal-stimulus conditions in the environmental assignment table, together
with the target-discovery rule and computational budgets, were frozen
before their outcomes were computed. The displayed cases in the main
article, \texttt{CD4\_APC} and \texttt{MAC\_BACTERIA}, were also selected
before the seven focal-context computations. No condition is replaced or omitted because of its outcome.

For completeness, the external-input sets are exactly the following.
All identifiers are literal repository coordinate names.
\begin{quote}\small
\textbf{CD4 (34 inputs):}
\path{v_APC}, \path{v_B7}, \path{v_CAV1_ACTIVATOR}, \path{v_CGC},
\path{v_ECM}, \path{v_GP130}, \path{v_GalphaS_L}, \path{v_Galpha_QL},
\path{v_Galpha_iL}, \path{v_IFNB_e}, \path{v_IFNGR1}, \path{v_IFNGR2},
\path{v_IFNG_e}, \path{v_IL10RA}, \path{v_IL10RB}, \path{v_IL10_e},
\path{v_IL12_e}, \path{v_IL15RA}, \path{v_IL15_e}, \path{v_IL18_e},
\path{v_IL21_e}, \path{v_IL22_e}, \path{v_IL23_e}, \path{v_IL27RA},
\path{v_IL27_e}, \path{v_IL2RB}, \path{v_IL2_e}, \path{v_IL4_e},
\path{v_IL6RA}, \path{v_IL6_e}, \path{v_IL9_e}, \path{v_MHC_II},
\path{v_TGFB_e}, \path{v_alpha_13L}.

\textbf{Macrophage (19 inputs):}
\path{v_Bacteria}, \path{v_CASP3_Activator}, \path{v_CASP6_Activator},
\path{v_CASP7_Activator}, \path{v_CD40LG}, \path{v_External_Activator},
\path{v_FASLG}, \path{v_IFNA}, \path{v_IFNB1}, \path{v_IFNG},
\path{v_ILIB}, \path{v_IRF1_IRF1_Activator}, \path{v_IRF2_Activator},
\path{v_Mitochondrial_Activator}, \path{v_RELA_NFKB1_Activator},
\path{v_TNF}, \path{v_TNFSF10}, \path{v_TNFSF13B}, \path{v_Virus}.
\end{quote}

% Complete table tab:largeinputs: biological_tables.tex


The CD4 focal inputs follow the subjects of Conroy et al.'s
Figure~2: APC, ECM, G$\alpha$q ligand and IL2~\cite{Conroy2014}.
The MHC-II coordinate in \texttt{CD4\_APC} and the CGC/IL2RB coordinates
in \texttt{CD4\_IL2} are declared cofactor conditions required by the
corresponding released receptor formulas. IL2RA remains dynamic. This
cofactor choice is not neutral: the released SHP2 rule also contains an
IL2RB term without ligand gating. The auxiliary CAV1 activator is zero
in all four focal conditions; this does not itself clamp the CAV1
scaffold to zero. These complete Boolean assignments are newly specified
computational environments, not reconstructions of the activity ranges
or stochastic environments in the source study. As the cofactor
backgrounds differ, cross-condition comparisons are not interpreted as
isolated causal effects of switching only one input.

For macrophages, IFNG and FASLG are separately traced in Raza et al.'s
Figure~6, and the Bacteria condition follows the pathogen-sensing scope
of their integrated pathway diagram~\cite{Raza2008}. Every focal condition
sets the auxiliary \path{v_External_Activator} to one. This repository
coordinate supplies such rows as IFNGR1, IFNGR2, JAK2, FAS and ProCASP8;
it is not a named biological ligand or an experimental concentration.
The Bacteria input feeds multiple pathogen-associated coordinates,
including LPS, Flagellin and diacyl/triacyl lipopeptides, and is not a
pure-LPS condition. Direct caspase, mitochondrial, IRF and RELA activators
remain zero in these three focal conditions. No healthy, diseased or
apoptotic phenotype is assumed from these assignments.

For each case, target discovery starts with every regulated coordinate
zero and the input coordinates at their recorded values. The model is
iterated synchronously for at most 2000 updates; only a full vector
satisfying every completed source formula is accepted as a target. A
repeated nontrivial orbit or an exhausted budget is retained as a
\emph{target not found by the declared procedure} status, not as a proof
that no fixed point exists. This deterministic choice of target is
separate from the subsequent all-initial-state control guarantee. the target-census section extends the analysis to every fixed target in the same frozen environments.
The node order is the lexicographically sorted BNET left-hand-side
names followed by the lexicographically sorted omitted input names.
The metadata's variable-name list is not substituted for the regulated
rows, because it omits some output coordinates. For each accepted case,
\path{results/CASE/certificate.json} records \texttt{node\_order}, the
full vector in \texttt{discovery.target}, the complete
\texttt{target\_by\_name} mapping and the fixed-point check. The run
manifest hashes each complete scientific certificate.


\subsection{Exact optimization and certificate scope}\label{sec:largecertificates}
For an environmental input set $B$ at its recorded values, we seek the
minimum additional target-pin set among all regulated coordinates.
Outcome coordinates are eligible under the same unit-cost rule as other
regulated coordinates. In particular, a mathematically optimal set may
contain an output such as proliferation; no gene-only or experimentally
implementable intervention restriction is inferred. Counting input
clamps gives the total cost $|B|+c_B$, with the main article's notation.
This identity follows
because every input row is an identity and hence must be pinned for
convergence from all full-model initial states. The additional cost
$c_B$ instead treats the complete environmental baseline as
fixed without cost.

The support procedure uses physical variables and exact local truth
tables after target-relative XOR relabeling. Its terminal cube supplies
a finite residual model. Exact global success on that residual is
equivalent to global success in the original completed model, by the
reduction result in the main article. Search therefore ranges over
subsets of the residual rather than all regulated coordinates. An
additional pin outside the residual has no effect on the restricted
map and cannot occur in a minimum additional set when it can be removed
without losing success. Nonempty residual support alone is not treated
as an unsuccessful control outcome.

\paragraph{Frozen search budgets and stopping rules.}
Each context has an analysis-time limit of 120 seconds, a residual
state-space limit of 18 coordinates, a search limit of six additional
pins, and a work cap of $30\,000\,000$. One work unit means one original
residual successor evaluation or one compatible residual state in a
candidate pin-set classification. The implementation reserves a whole
candidate's state count before evaluating it; reserved work is therefore
a conservative accounting in an interrupted evaluation. These counts
are distinct from source-formula parsing, target-discovery updates and
support computations. Repeated occurrences of one physical variable
share a single truth-table argument. Local tables compile every Boolean
corner: 3176 corners for the completed CD4
map and 3780 for the completed macrophage map, with maximum physical
arity ten in either model.

For an enumerable residual of size $r$, the original successor table
contains $2^r$ states. Candidate sets are visited in increasing
cardinality and lexicographic coordinate order. Each candidate is checked
on all $2^{r-k}$ compatible states, including every directed cycle of
its deterministic functional graph. Every smaller cardinality must be
excluded before reporting an exact optimum. At the first successful
cardinality, search continues to enumerate alternative minimum sets
until that level is complete or a budget is exhausted. Thus an exact
minimum cost can be certified even when the list of all minimizers is
incomplete. If a residual exceeds the enumeration cap, the declared
fallback is an exact pure-copy certificate when applicable; otherwise
the case retains bounds and a residual-size-cap status. A budget or
cardinality cap never becomes a claim of infeasibility. Full state
spaces compatible with the respective input baselines have sizes
$2^{154}$ and $2^{302}$ and are not enumerated.

\subsection{Original selected-target certificates}\label{sec:largeresults}
the selected-target outcome table retains all eleven conditions. Ten target
vectors are found and independently checked as full fixed points.
All ten associated minimum costs are exact. In \texttt{MAC\_BACTERIA},
the declared initial orbit enters a period-four cycle at update 11 and
repeats at update 15, so no fixed target is selected. The subsequent exhaustive check in the target-census section proves that this baseline-clamped map has no fixed point. The recorded orbit remains a separate dynamical observation.
Main-article The macrophage cycle panel in \path{figures/fig5_larger_network_support.pdf} displays the 19 coordinates varying around this
verified four-cycle, at updates 11--15. Its final column repeats the
first. The remaining 302 coordinates are constant on the cycle,
comprising 283 regulated coordinates and all 19 inputs. The full orbit,
the varying-coordinate projection and all constant values are supplied
as plotting data. Panel B shows actual Boolean states, whereas panel A
shows coordinate-cube support cardinalities.

% Complete table tab:largeoutcomes: biological_tables.tex


The four convergent macrophage cases require no additional intervention:
their full support becomes empty after the recorded baseline-input
assignment. The six CD4 residuals require exact finite-state analysis.
The complete optimum families for \texttt{CD4\_OFF}, \texttt{CD4\_ECM}
and \texttt{CD4\_GALPHAQ} are
\[
 \{\texttt{v\_Dec2},\texttt{v\_TBET},\texttt{v\_proliferation}\},
 \qquad
 \{\texttt{v\_GATA3},\texttt{v\_TBET},\texttt{v\_proliferation}\},
\]
with every listed coordinate pinned to zero.
For \texttt{CD4\_IL2}, only the second of these sets is optimal.
For \texttt{CD4\_ON}, the two optimum sets are
\[
 \{\texttt{v\_Dec2},\texttt{v\_FOXP3},\texttt{v\_TBET}\},
 \qquad
 \{\texttt{v\_FOXP3},\texttt{v\_GATA3},\texttt{v\_TBET}\},
\]
where TBET is pinned to one and the other listed coordinates to zero.
These are statements about model coordinates and the declared target,
not proposed molecular intervention protocols.

For \texttt{CD4\_APC}, all
$\sum_{j=0}^{3}\binom{17}{j}=834$ smaller candidate sets are excluded.
Two feasible four-coordinate sets are certified before the work cap:
\[
\begin{gathered}
 \{\texttt{v\_Dec2},\texttt{v\_FOXP3},\texttt{v\_TBET},
       \texttt{v\_proliferation}\},\\
 \{\texttt{v\_FOXP3},\texttt{v\_GATA3},\texttt{v\_TBET},
       \texttt{v\_proliferation}\},
\end{gathered}
\]
all pinned to zero. These establish the exact additional cost four.
The run checks 2424 candidate sets in total and reserves
$29\,999\,104$ work units; the next candidate would exceed the cap.
The original whole-residual run did not exhaust the four-node minimizers. The SCC analysis in the target-census section now proves that the two displayed controls are the complete optimum family. The original case record is retained with its recorded scope, including its search status and the witness used for its transient bound.

\paragraph{Measured computational resources.}
the resource table reports the original selected-target run, with one isolated worker process per
context under Python 3.12.14 on Linux x86\_64. Analysis wall time and
CPU time cover model parsing, target discovery, support reduction,
residual search and construction of the certificate object; they exclude
interpreter startup, imports and file serialization. Peak resident-set
size includes interpreter/import overhead. Full-process wall time and
unrounded CPU time are retained in \path{resources.json}. The cap applies
to the analysis clock. The measurements describe this run and are not
hardware-independent performance guarantees.

% Complete table tab:largeresources: biological_tables.tex




\section{All fixed targets and functional feedback components}\label{sec:alltargets}
We retain the eleven environmental assignments and original model files,
then enumerate all fixed points rather than selecting only the zero-orbit
target. This extension was performed after the original selected-target
study. Its settings and results are recorded separately in
\path{results/all_targets.json}; the original protocol and scientific
certificates remain the same.

\subsection{Enumeration, localization and verification}
Start with the complete environmental assignment. If a regulated rule
has one constant value for every assignment of its remaining free inputs,
that coordinate is forced to that value in every fixed point. Iterate
this rule until no further coordinate is forced. Every fixed point lies
in the resulting cube, so testing all assignments of its free coordinates
against every source equation yields the complete fixed-point set.
Across the eleven contexts, this examines $805380$ assignments, split as
$281088$ for CD4 and $524292$ for macrophage. In the
macrophage BACTERIA context, $19$ coordinates remain free and none of
their $2^{19}=524288$ assignments is a fixed point. The other four
macrophage contexts each have one fixed point and empty target residual.

For each of the $46$ CD4 and four macrophage targets, compute the exact
physical support residual. On its cube, a parent is retained in a row
precisely when some Boolean-difference test for that parent is nonzero.
Testing all local assignments includes cooperative influences with zero
singleton sensitivity. The resulting exact functional graph contains
$478$ SCCs across all targets, including acyclic singleton components.
The largest SCC has eleven coordinates, so each component control problem
has at most $2^{11}$ local states. The main article's Theorem~2.5
applies to this graph. An SCC's isolated truth table is formed by a
further substitution of all exterior coordinates to their target values;
this substitution can remove dependencies that are present on the full
residual cube. Both tables are retained in the machine-readable record.

The solver enumerates candidate pin sets by cardinality and tests all
compatible local states. It then enumerates every local minimum set and
checks a combined witness on every compatible residual state. All $50$
costs and all local minimum families are exact. Global minimum families
are the Cartesian products of the local families, by SCC localization.
In particular, the selected APC target has exactly the two optimal sets
listed in the larger-model protocol section.3.

The declared guards allow at most $24$ free coordinates for fixed-point
enumeration, in chunks of $8192$ assignments. Component optimization uses
size and pin-count caps of $18$ and a cap of $10^8$ work units per complex
component. A work unit has the SCC solver's definition, one local-row
evaluation or controlled successor-edge traversal. Enumeration of all
local minimizers and of restricted eligible sets has a separate component
size cap of $12$. No guard is exhausted in this study. These settings
describe this complete finite analysis and do not assert tractability
for arbitrary models.

An independent recursive parser and an exhaustive Boolean-mask enumerator
reproduce the fixed-point set from the original source rules. Separate
support, functional-dependence, mutual-reachability and reverse-reachability
calculations verify all $478$ SCC records and all local minimum families.
The $50$ global witnesses are checked over $281092$ compatible residual
states in total. The seven feasible witnesses with proliferation forbidden
are checked over $49152$ compatible residual states. These counts include
repeated states across different target and intervention problems.

\subsection{Complete target table}
the complete CD4 target table lists every CD4 fixed target. A marker
signature records the literal values of TBET, GATA3, FOXP3, RORGT,
proliferation and STAT5\_HIGH in that order. These signatures describe
model coordinates and are not assigned unique lineage or health labels.
IDs follow the deterministic enumeration order in the certificate;
an asterisk identifies the original zero-orbit target. Full ordered
$188$-bit target vectors, all named coordinate values and additional
marker values are provided in the computational record.

In the witness column, $D=\mathrm{Dec2}$, $G=\mathrm{GATA3}$,
$F=\mathrm{FOXP3}$, $H=\mathrm{IL4RA\_HIGH}$, $R=\mathrm{RORGT}$,
$T=\mathrm{TBET}$ and $P=\mathrm{proliferation}$.
A subscript gives the permanent target value. The column $M$ counts
all minimum additional pin sets, obtained from the local families.
The restricted cost $c_B^{\neg P}$ forbids proliferation as an intervention
coordinate; $\infty$ denotes proved infeasibility.
% Complete table tab:allcd4targets: biological_tables.tex


\subsection{Eligibility and the meaning of the component decomposition}
For a fixed eligible set $E$, the pin-set equivalence in the main article's
Theorem~2.5 gives the same additive decomposition with
\[
 \kappa_E(C)=\min\{|J|:J\subseteq E\cap C,\ J\text{ stabilizes }F_C\},
\]
where the minimum is $\infty$ if no eligible set succeeds. Every successful
global set restricts to a successful local set; conversely the union of
local successful sets succeeds. Thus infeasibility of one isolated
component proves infeasibility of the full target problem. This is a
restriction of the allowable pin sets, not an assumption that successful
sets are monotone under adding pins.

Exactly $39$ CD4 targets have an isolated proliferation identity component.
All have target STAT5\_HIGH equal to zero. If proliferation is ineligible,
the local eligible set is empty and the opposite value is a competing
fixed point. This includes target proliferation equal to one. In the six
ON targets, support reduction removes proliferation. In target IL2 05,
it remains in the residual, but target STAT5\_HIGH is one and the
isolated proliferation rule is constant at its target value one. These
seven problems retain their original optimum when proliferation is
excluded. Hence infeasibility,
rather than an increased finite cost, is the appropriate entry in
the complete CD4 target table.

The functional SCC partition need not consist only of positive-cost
components. For the selected APC target, the RORGT singleton has a
self-dependency somewhere on the full residual cube, but its isolated
rule is zero. It therefore costs zero. Similarly, the four-coordinate
FOXP3 component has isolated rules
\[
 f^+=f\vee s,\qquad s^+=r,\qquad t^+=0,\qquad r^+=t,
\]
in the order FOXP3, SMAD3, TGFB, TGFBR. Its external target substitution
breaks the TGFB feedback, but FOXP3 retains a self-memory and uniquely
requires a pin. This explains why a count of cyclic SCCs alone is
insufficient outside the copy-certificate hypotheses.

\section{Fair asynchronous updates for two selected targets}\label{sec:fairupdates}
% The general residual-fairness proposition and full proof are now in the
% main article, Proposition 7.1, label prop:fair-coordinate-reduction.
% They are intentionally not duplicated in this supporting excerpt.


\paragraph{A selected-target asynchronous check.}
We additionally test the archived CD4 OFF and APC targets under these
coordinate-fair updates, retaining their input assignments and intervention
convention. Direct graph enumeration certifies the pin sets
\[
\begin{aligned}
J_{\mathrm{OFF}}&=\{\mathrm{Dec2},\mathrm{TBET},\mathrm{proliferation}\},\\
J_{\mathrm{APC}}&=J_{\mathrm{OFF}}\cup\{\mathrm{FOXP3}\},
\end{aligned}
\]
and the two variants obtained by replacing Dec2 with GATA3. After applying
the pins to the baseline residual, each OFF graph has $64$ states and
$192$ state-changing edges; each APC graph has $8192$ states and $51200$
edges. All four graphs are acyclic with zero as their only sink.
The main article's residual-reduction proposition under coordinate fairness
(label \texttt{prop:fair-coordinate-reduction}) therefore proves convergence from every
compatible full-network initial state under every coordinate-fair schedule.

The lower bounds are verified independently of synchronous optimization.
In the OFF residual, the disjoint functional SCCs
$\{\mathrm{Dec2},\mathrm{GATA3}\}$, $\{\mathrm{TBET}\}$, and
$\{\mathrm{proliferation}\}$ each admit a nonzero isolated fixed point
when unpinned. In the displayed node orders these witnesses are $(1,1)$,
$1$, and $1$. APC has the same three obstructions and the additional SCC
$\{\mathrm{FOXP3},\mathrm{SMAD3},\mathrm{TGFB},\mathrm{TGFBR}\}$,
with isolated fixed point $(1,0,0,0)$. All listed target coordinates are
zero. Initializing the strict ancestor components at their target values
preserves each such local fixed point under arbitrary target pins outside
that SCC, by the necessity argument of
the main article's Theorem~2.5. This argument also holds for
one-coordinate updates. Every successful intervention must therefore
intersect all three, respectively four, components. Together with the
four graph certificates, this establishes exact additional optima of
three for OFF and four for APC under coordinate-fair updating.

The acyclicity check is stronger than requiring a unique terminal SCC.
For example, $f_1=x_1x_2$, $f_2=x_1(1-x_2)$ has $(0,0)$ as its only
terminal asynchronous component, but repeatedly updating coordinates
$2,1,2$ from $(1,0)$ gives
$(1,0)\to(1,1)\to(1,1)\to(1,0)$.
This schedule is coordinate-fair and never reaches the target.
The CD4 certificates exclude this possibility by checking the entire
state-changing residual graph. They concern these selected targets and
update semantics; the random-ensemble results retain their synchronous
formulation.

\paragraph{Eligibility and the proliferation memory rule.}
The supplied CD4 rule is
\[
\mathrm{proliferation}^{+}
=\mathrm{STAT5\_HIGH}\vee\mathrm{proliferation}.
\]
For a target with proliferation equal to zero, leaving this coordinate
unpinned makes the full-state target unattainable from every compatible
initial state: an initial proliferation value of one remains one at
every subsequent update, independently of all other interventions.
The argument applies to synchronous and coordinate-fair asynchronous
updates. Excluding proliferation from eligibility therefore makes the
selected OFF, APC, ECM, IL2, and GALPHAQ target problems infeasible.
It does not merely increase their finite intervention cost.

Proliferation already contributes one mandatory additional pin in OFF.
The extra pin in APC is associated with the FOXP3 component, rather than
with a newly appearing proliferation requirement. These conclusions
describe the consequence of the encoded absorbing-memory rule and the
chosen full-state targets. A target with proliferation equal to one
requires its own eligibility analysis: when its isolated residual rule is
the identity, a missing eligible pin again leaves a competing fixed point.
The absorbing-one argument alone does not settle that case. Replacing a
full-state target by a phenotype or changing the eligible coordinates
defines a different control problem.


\section{Reproduction and validation}
\label{bio:reproduction}
The biological workflow is \path{workflows/biological/}; paths in this
reference are relative to that directory. It contains the rules, targets,
certificates and reproduction inputs. The separate \path{workflows/scc/}
workflow provides the SCC-localized solver and independent finite-state checks.
Use the repository-level runner described in \path{docs/REPRODUCTION.md} for
isolated execution. Source lineage and recorded commitments are distinguished
from current representation hashes under \path{provenance/}.


The two smaller benchmarks are reproduced by the \path{biological/}
package, independently of the random-network simulations. Its entry
point, relative to that directory, is \path{src/reproduce.py}; the README gives the command and
output-directory convention. Model extraction uses a restricted syntax
tree evaluator and exact arithmetic. Complete JSON tables, Boolean BNET
exports, support traces, witness trajectories, and finite-state counts
are provided. Source checksums identify the particular input versions.
The engine checks 256 exhaustive two-node map--target pairs, 72 deterministically generated three- to five-node models, and a repeated-variable case. In total, 329 models, 2052 support cubes, 7713 baseline/additional-pin combinations, and 1680 pure-copy certificates are checked. Independent physical-state enumeration verifies the support reduction, intervention feasibility and pure-copy optima. A nonmonotone pinning example checks that physical success is never presumed to persist under arbitrary additional pins. These finite tests validate the implementation; the general result follows from the proof in the main article.


The larger models have a separate \path{larger_networks/} package.
From that directory, the following commands rerun the frozen experiment
into a fresh output directory and independently verify its certificates:
\begin{quote}\small
\begin{verbatim}
python src/benchmark.py --protocol protocol/contexts.json \
  --output results_reproduced
python verification/verify_results.py \
  --protocol protocol/contexts.json --results results_reproduced \
  --output verification/reproduced_audit.json
python src/test_runner.py
\end{verbatim}
\end{quote}
The runner refuses to overwrite an existing output directory. The
independent verifier reads the frozen BNET with a separate recursive
parser, evaluates all local truth corners, recomputes target orbits and
physical cube propagation, and checks residual control through reverse
reachability. It imports neither the production parser nor its control
engine. Every stored failure-cycle edge is lifted to the original full
Boolean map with outside-residual coordinates at target. It also checks
candidate cardinality coverage before accepting a lower bound and
separately verifies each recorded successful set. The additional small
search tests cover all 256 two-coordinate zero-fixed-map/pin combinations,
nonmonotone pinning, work/time/cardinality caps, the pure-copy fallback
and cyclic target discovery.

The packaged independent audit verifies all eleven case records, including
the bacterial discovery cycle. The all-target extension separately proves absence of a fixed point in that context. It reclassifies 4210 candidate pin sets
over $48\,038\,912$ compatible candidate-state pairs, summed across
sets and contexts, recomputes 281088
residual successors, and checks 6511 failure-cycle edges after lifting
them to their original full networks. The report is
\path{larger_networks/verification/independent_results.json}. These
verification counts are reported separately from the production work
budgets and resource measurements.

\begin{samepage}
The fixed protocol, sources, certificates and hashes identify the
reproducible results. Runtime and memory use may vary, and an earlier wall
cutoff can limit search progress. Implementation checks and exact model
convergence do not experimentally validate the full target states.\par
\end{samepage}


\begin{samepage}
\paragraph{All-target and fair-update reproduction.}
From an isolated copy of the biological workflow, run
\begin{quote}\small\begin{verbatim}
python src/run_all_targets.py --output results_reproduced
python verification/verify_all_targets.py \
  --production results_reproduced/all_targets.json
python verification/async_memory_check.py
\end{verbatim}\end{quote}
\end{samepage}
The first command regenerates the target census and SCC certificates.
The independent verifier recomputes the expected results from source and
compares them with the regenerated \path{results_reproduced/all_targets.json}.
The third command
uses its own parser and checks the selected-target fair-update graphs
and local fixed-point obstructions. Deterministic scientific records are
separated from timing information. The figure exporter reads these exact
counts and separately checks the displayed isolated rules against the
frozen BNET expressions.


\begin{thebibliography}{9}\small\setlength{\itemsep}{0pt}
\bibitem{Zhu2024} S. Zhu, J. Cao, L. Lin, J. Lam and S.-i. Azuma,
Toward stabilizable large-scale Boolean networks by controlling the minimal set of nodes,
\emph{IEEE Transactions on Automatic Control} \textbf{69} (2024), 174--188.
\href{https://doi.org/10.1109/TAC.2023.3269321}{doi:10.1109/TAC.2023.3269321}.
\bibitem{Parmer2022} T. Parmer, L. M. Rocha and F. Radicchi,
Influence maximization in Boolean networks,
\emph{Nature Communications} \textbf{13} (2022), 3457.
\href{https://doi.org/10.1038/s41467-022-31066-0}{doi:10.1038/s41467-022-31066-0}.
\bibitem{Li2006} S. Li, S. M. Assmann and R. Albert,
Predicting essential components of signal transduction networks: A dynamic model of guard cell abscisic acid signaling,
\emph{PLOS Biology} \textbf{4} (2006), e312.
\href{https://doi.org/10.1371/journal.pbio.0040312}{doi:10.1371/journal.pbio.0040312}.
\bibitem{Conroy2014} B. D. Conroy, T. A. Herek, T. D. Shew, M. Latner,
J. J. Larson, L. Allen, P. H. Davis, T. Helikar and C. E. Cutucache,
Design, assessment, and in vivo evaluation of a computational model
illustrating the role of CAV1 in CD4$^+$ T-lymphocytes,
\emph{Frontiers in Immunology} \textbf{5} (2014), 599.
\href{https://doi.org/10.3389/fimmu.2014.00599}{doi:10.3389/fimmu.2014.00599}.
\bibitem{Raza2008} S. Raza, K. A. Robertson, P. A. Lacaze, D. Page,
A. J. Enright, P. Ghazal and T. C. Freeman,
A logic-based diagram of signalling pathways central to macrophage activation,
\emph{BMC Systems Biology} \textbf{2} (2008), 36.
\href{https://doi.org/10.1186/1752-0509-2-36}{doi:10.1186/1752-0509-2-36}.
\bibitem{Helikar2012} T. Helikar, B. Kowal, S. McClenathan, M. Bruckner,
T. Rowley, A. Madrahimov, B. Wicks, M. Shrestha, K. Limbu and J. A. Rogers,
The Cell Collective: Toward an open and collaborative approach to systems biology,
\emph{BMC Systems Biology} \textbf{6} (2012), 96.
\bibitem{Pastva2023} S. Pastva, D. \v{S}afr\'{a}nek, N. Bene\v{s}, L. Brim and T. Henzinger,
Repository of logically consistent real-world Boolean network models,
\emph{bioRxiv} (2023), 2023.06.12.544361.
\href{https://doi.org/10.1101/2023.06.12.544361}{doi:10.1101/2023.06.12.544361}.
\end{thebibliography}
